\section{Conclusions}
\label{sec:conclusions}

We have developed PULSUS, a direct frequency-domain framework for
multidimensional spectroscopy of open quantum systems that retains
finite pulse duration and, when required, the complete light--matter
interaction beyond the rotating-wave approximation.
The method combines Liouville-space resolvents with operator-valued
Gaussian pulse dressing, allowing the excitation and detection
frequency coordinates to be evaluated directly without first
constructing a time-domain interferogram.
The same formulation generates a controlled hierarchy ranging from the
full finite-pulse response to the conventional impulsive-RWA limit.

The finite-pulse implementation was validated against the analytical
Gaussian-pulse benchmark of Smallwood \textit{et al.}, yielding a
relative discrepancy of $7.33\times10^{-6}$ at
$T=2\pi\sigma$ and approaching the analytical result still more closely
as the pulse separation is increased.
Independent Cartesian-grid and arbitrary-coordinate implementations
agree at numerical precision.
The latter permits spectra to be evaluated only at requested
$(\omega_1,\omega_3)$ pairs and therefore provides a natural interface
for targeted spectral sampling.
For the coupled-dimer benchmark considered here, redistributing the
sampling budget toward the spectrally active region reduced the number
of evaluated coordinates required for approximately $10\%$
reconstruction error by a factor of about $6.2$, with a corresponding
wall-clock reduction of about $1.8$.
These numerical gains are model and reconstruction dependent, but
demonstrate the utility of treating arbitrary spectral coordinates as
first-class computational objects.

The controlled hierarchy
$A\rightarrow B\rightarrow C\rightarrow D$
also makes it possible to separate physically distinct sources of
error in the conventional impulsive-RWA description.
For the coupled dimer, the $A-B$ comparison isolates non-rotating-wave
contributions, $B-C$ isolates molecular evolution occurring during the
pulse envelopes, and $C-D$ isolates the remaining finite-bandwidth
effect.
The resulting regime map shows that these errors dominate in different
pulse-duration ranges.
Very short pulses suppress in-pulse molecular evolution but can expose
counter-rotating pathways, whereas long pulses enhance both molecular
evolution during the interaction and finite-bandwidth effects.
Consequently, the total impulsive-RWA error is nonmonotonic and is
minimized in an intermediate regime rather than in the shortest-pulse
limit.

For the nonrephasing response, the short-pulse breakdown of the RWA
can be traced to specific field-sign sectors of the complete
interaction.
In the present model, the dominant discarded contribution has
signature $(+,+,-)$.
This contribution vanishes in the uncoupled limit and grows with
electronic coupling, reaching approximately $45\%$ of the conventional
$(+,-,+)$ nonrephasing contribution at
$J=40~\mathrm{meV}$ for a pulse intensity FWHM of
$0.10~\mathrm{fs}$.
The result illustrates that RWA validity is controlled not only by the
pulse duration relative to the carrier period, but also by the
molecular couplings that determine which counter-rotating pathways are
activated.

The present formulation assumes a time-independent field-free
Liouvillian and temporally separated pulses, with the finite-pulse
results developed here specialized to Gaussian envelopes.
Within these conditions, PULSUS provides a direct route from an
open-system molecular model to finite-pulse multidimensional spectra
while retaining explicit control over the approximations normally
introduced in constructing the signal.
This combination of direct spectral evaluation, selective frequency
sampling, and approximation-resolved analysis provides a practical
framework for determining when the impulsive rotating-wave description
is reliable and when the neglected finite-pulse or non-rotating-wave
physics must be retained.